
\documentclass{article}
\usepackage{colortbl}
\usepackage{makecell}
\usepackage{multirow}
\usepackage{supertabular}

\begin{document}

\newcounter{utterance}

\centering \large Interaction Transcript for game `function\_detective', experiment `numbers\_passive\_examples', episode 1 with Qwen3.6{-}35B{-}A3B{-}FP8{-}without{-}reasoning.
\vspace{24pt}

{ \footnotesize  \setcounter{utterance}{1}
\setlength{\tabcolsep}{0pt}
\begin{supertabular}{c@{$\;$}|p{.15\linewidth}@{}p{.15\linewidth}p{.15\linewidth}p{.15\linewidth}p{.15\linewidth}p{.15\linewidth}}
    \# & \multicolumn{2}{c}{Player} && \multicolumn{2}{c}{Game Master} \\
    \hline

    \theutterance \stepcounter{utterance}  
    & & & \multicolumn{4}{p{0.6\linewidth}}{
        \cellcolor[rgb]{0.9,0.9,0.9}{
            \makecell[{{p{\linewidth}}}]{
                \texttt{\tiny{[P1$\langle$GM]}}
                \texttt{I have a hidden Python function for you to figure out.} \\
\\ 
\texttt{Your job is to infer what it does and then submit a correct Python implementation.} \\
\\ 
\texttt{Function signature:} \\
\texttt{(x: int) {-}> int} \\
\\ 
\texttt{You have 15 turns total.} \\
\\ 
\texttt{I control which valid input{-}output examples you see. You can request them one at a time.} \\
\\ 
\texttt{You must make your best SOLVE attempt no later than the second{-}to{-}last turn.} \\
\texttt{Do not wait until the very end.} \\
\\ 
\texttt{You can use these actions:} \\
\texttt{{-} JSON format described at the end of the prompt} \\
\texttt{{-} SOLVE: ```python ... ```} \\
\\ 
\texttt{Do not add import statements in your guess.} \\
\texttt{Anything you might need is already available, so just write the function body in hypothesis and return the answer.} \\
\\ 
\texttt{Preloaded examples:} \\
\texttt{1. ({-}10, {-}11)} \\
\\ 
\texttt{Scoring:} \\
\texttt{{-} correctness is required for any non{-}zero score} \\
\texttt{{-} in interactive modes, solving with fewer information requests is better} \\
\\ 
\texttt{Input formatting reminders:} \\
\texttt{{-} strings must be quoted} \\
\texttt{{-} booleans must be True/False} \\
\texttt{{-} use commas between multiple values} \\
\texttt{{-} do not write any extra action tags in the same message} \\
\\ 
\texttt{Format example:} \\
\texttt{You write:} \\
\texttt{\{'current\_mode': <probing/confirming/don't know; choose whichever best describes your current state in the task>, 'input\_rationale': <brief description of why you chose this hypothesis>, 'hypothesis': <either the exact string "unknown" if you have no guess yet OR your current best guess as runnable Python matching the given signature, with a one{-}line docstring stating what it does. Never write placeholder code, if you cannot commit to a guess, use "unknown". Match the given signature.} \\
\texttt{e.g.} \\
\texttt{def f(x):} \\
\texttt{	"""Returns the absolute value of x."""} \\
\texttt{	return x if x >= 0 else {-}x} \\
\\ 
\texttt{Use "unknown" if you have no guess.>, 'confidence\_score': <score between 1{-}100. 100 means highest confidence in your hypothesis>\}} \\
\\ 
\texttt{Output only one valid json per turn.} \\
            }
        }
    }
     \\ \\

    \theutterance \stepcounter{utterance}  
    & \multicolumn{4}{p{0.6\linewidth}}{
        \cellcolor[rgb]{0.561,0.737,0.561}{
            \makecell[{{p{\linewidth}}}]{
                \texttt{\tiny{[P1$\rangle$GM]}}
                \texttt{\{'current\_mode': 'probing', 'input\_rationale': 'I have one example: f({-}10) = {-}11. This suggests a simple linear relationship like f(x) = x {-} 1. Let me test this hypothesis with another input to confirm.', 'hypothesis': 'unknown', 'confidence\_score': 20\}} \\
            }
        }
    }
    & & \\ \\

    \theutterance \stepcounter{utterance}  
    & & & \multicolumn{4}{p{0.6\linewidth}}{
        \cellcolor[rgb]{0.9,0.9,0.9}{
            \makecell[{{p{\linewidth}}}]{
                \texttt{\tiny{[P1$\langle$GM]}}
                \texttt{Game Violation: ParseError: Response is not valid JSON.} \\
\texttt{Output ONLY:} \\
\texttt{One JSON with confidence\_score, current\_mode, input\_rationale, and hypothesis keys ...} \\
\texttt{OR} \\
\texttt{SOLVE: ```python ... ```} \\
            }
        }
    }
     \\ \\

    \theutterance \stepcounter{utterance}  
    & \multicolumn{4}{p{0.6\linewidth}}{
        \cellcolor[rgb]{0.561,0.737,0.561}{
            \makecell[{{p{\linewidth}}}]{
                \texttt{\tiny{[P1$\rangle$GM]}}
                \texttt{\{'current\_mode': 'probing', 'input\_rationale': 'I have one example: f({-}10) = {-}11. This suggests a simple linear relationship like f(x) = x {-} 1. Let me test this hypothesis with another input to confirm.', 'hypothesis': 'unknown', 'confidence\_score': 20\}} \\
            }
        }
    }
    & & \\ \\

    \theutterance \stepcounter{utterance}  
    & & & \multicolumn{4}{p{0.6\linewidth}}{
        \cellcolor[rgb]{0.9,0.9,0.9}{
            \makecell[{{p{\linewidth}}}]{
                \texttt{\tiny{[P1$\langle$GM]}}
                \texttt{Game Violation: ParseError: Response is not valid JSON.} \\
\texttt{Output ONLY:} \\
\texttt{One JSON with confidence\_score, current\_mode, input\_rationale, and hypothesis keys ...} \\
\texttt{OR} \\
\texttt{SOLVE: ```python ... ```} \\
            }
        }
    }
     \\ \\

    \theutterance \stepcounter{utterance}  
    & \multicolumn{4}{p{0.6\linewidth}}{
        \cellcolor[rgb]{0.561,0.737,0.561}{
            \makecell[{{p{\linewidth}}}]{
                \texttt{\tiny{[P1$\rangle$GM]}}
                \texttt{\{"current\_mode": "probing", "input\_rationale": "I have one example: f({-}10) = {-}11. This suggests a simple linear relationship like f(x) = x {-} 1. Let me test this hypothesis with another input to confirm.", "hypothesis": "unknown", "confidence\_score": 20\}} \\
            }
        }
    }
    & & \\ \\

    \theutterance \stepcounter{utterance}  
    & & & \multicolumn{4}{p{0.6\linewidth}}{
        \cellcolor[rgb]{0.9,0.9,0.9}{
            \makecell[{{p{\linewidth}}}]{
                \texttt{\tiny{[P1$\langle$GM]}}
                \texttt{OUTPUT: (0, 0)} \\
            }
        }
    }
     \\ \\

    \theutterance \stepcounter{utterance}  
    & \multicolumn{4}{p{0.6\linewidth}}{
        \cellcolor[rgb]{0.561,0.737,0.561}{
            \makecell[{{p{\linewidth}}}]{
                \texttt{\tiny{[P1$\rangle$GM]}}
                \texttt{\{"current\_mode": "probing", "input\_rationale": "I have two examples: f({-}10) = {-}11 and f(0) = 0. The hypothesis f(x) = x {-} 1 is incorrect because f(0) should be {-}1, not 0. Let's try f(x) = x + c*x for some c. f(0)=0 holds. f({-}10) = {-}10 + c*({-}10) = {-}11 => {-}10(1+c) = {-}11 => 1+c = 1.1 => c = 0.1. So f(x) = 1.1*x? But return type is int. Maybe f(x) = x + x//10? Let's check: f({-}10) = {-}10 + ({-}10)//10 = {-}10 + ({-}1) = {-}11. f(0) = 0 + 0 = 0. This works. Let me test with another value to confirm.", "hypothesis": "unknown", "confidence\_score": 60\}} \\
            }
        }
    }
    & & \\ \\

    \theutterance \stepcounter{utterance}  
    & & & \multicolumn{4}{p{0.6\linewidth}}{
        \cellcolor[rgb]{0.9,0.9,0.9}{
            \makecell[{{p{\linewidth}}}]{
                \texttt{\tiny{[P1$\langle$GM]}}
                \texttt{OUTPUT: ({-}2, {-}4)} \\
            }
        }
    }
     \\ \\

    \theutterance \stepcounter{utterance}  
    & \multicolumn{4}{p{0.6\linewidth}}{
        \cellcolor[rgb]{0.561,0.737,0.561}{
            \makecell[{{p{\linewidth}}}]{
                \texttt{\tiny{[P1$\rangle$GM]}}
                \texttt{\{"current\_mode": "probing", "input\_rationale": "I have three examples: f({-}10) = {-}11, f(0) = 0, f({-}2) = {-}4. Let's check the hypothesis f(x) = x + x//10. For x={-}2: {-}2 + ({-}2)//10 = {-}2 + ({-}1) = {-}3. But the output is {-}4. So this hypothesis is incorrect. Let's try f(x) = 2*x. f({-}10) = {-}20, not {-}11. Let's try f(x) = x {-} 1 for x < 0 and f(x) = x for x >= 0? f({-}10) = {-}11, f(0) = 0, f({-}2) = {-}3. But f({-}2) is {-}4. So that's not it either. Let's try f(x) = x + floor(x/10) but with different rounding? Or maybe f(x) = x {-} ceil(|x|/10)? For x={-}10: {-}10 {-} 1 = {-}11. For x=0: 0 {-} 0 = 0. For x={-}2: {-}2 {-} ceil(2/10) = {-}2 {-} 1 = {-}3. Still not {-}4. Let's try f(x) = x {-}} \\
            }
        }
    }
    & & \\ \\

    \theutterance \stepcounter{utterance}  
    & & & \multicolumn{4}{p{0.6\linewidth}}{
        \cellcolor[rgb]{0.9,0.9,0.9}{
            \makecell[{{p{\linewidth}}}]{
                \texttt{\tiny{[P1$\langle$GM]}}
                \texttt{Game Violation: ParseError: Response is not valid JSON.} \\
\texttt{Output ONLY:} \\
\texttt{One JSON with confidence\_score, current\_mode, input\_rationale, and hypothesis keys ...} \\
\texttt{OR} \\
\texttt{SOLVE: ```python ... ```} \\
            }
        }
    }
     \\ \\

    \theutterance \stepcounter{utterance}  
    & \multicolumn{4}{p{0.6\linewidth}}{
        \cellcolor[rgb]{0.561,0.737,0.561}{
            \makecell[{{p{\linewidth}}}]{
                \texttt{\tiny{[P1$\rangle$GM]}}
                \texttt{\{"current\_mode": "probing", "input\_rationale": "I have three examples: f({-}10) = {-}11, f(0) = 0, f({-}2) = {-}4. The hypothesis f(x) = x + x//10 failed for x={-}2. Let's try f(x) = 2*x for negative numbers? f({-}10) = {-}20 (no). Let's try f(x) = x {-} 1 for x <= {-}10, and f(x) = 2*x for {-}10 < x < 0? f({-}10) = {-}11, f({-}2) = {-}4, f(0) = 0. This fits all three. Let's test with another negative number like {-}5 to confirm.", "hypothesis": "unknown", "confidence\_score": 50\}} \\
            }
        }
    }
    & & \\ \\

    \theutterance \stepcounter{utterance}  
    & & & \multicolumn{4}{p{0.6\linewidth}}{
        \cellcolor[rgb]{0.9,0.9,0.9}{
            \makecell[{{p{\linewidth}}}]{
                \texttt{\tiny{[P1$\langle$GM]}}
                \texttt{OUTPUT: ({-}99, {-}198)} \\
            }
        }
    }
     \\ \\

    \theutterance \stepcounter{utterance}  
    & \multicolumn{4}{p{0.6\linewidth}}{
        \cellcolor[rgb]{0.561,0.737,0.561}{
            \makecell[{{p{\linewidth}}}]{
                \texttt{\tiny{[P1$\rangle$GM]}}
                \texttt{\{"current\_mode": "probing", "input\_rationale": "I have four examples: f({-}10) = {-}11, f(0) = 0, f({-}2) = {-}4, f({-}99) = {-}198. The hypothesis f(x) = 2*x for x < 0 and f(x) = x for x >= 0 fits all points: f({-}10) = {-}20 (wait, {-}20 != {-}11). So that hypothesis is wrong. Let's re{-}examine. f({-}10)={-}11, f({-}99)={-}198. {-}198 is exactly 2 * {-}99. {-}11 is not 2 * {-}10. So it's not simply 2*x. Let's look at f({-}10)={-}11. Maybe f(x) = x + floor(x/10) was close but failed for {-}2. What if it's f(x) = x + floor(x/10) for x <= {-}10, and something else for {-}10 < x < 0? Or maybe a different formula. Let's try f(x) = x {-} 1 for x <= {-}10? f({-}10) = {-}11. f({-}99) = {-}100 (no, it's {-}198). So for large negatives, it} \\
            }
        }
    }
    & & \\ \\

    \theutterance \stepcounter{utterance}  
    & & & \multicolumn{4}{p{0.6\linewidth}}{
        \cellcolor[rgb]{0.9,0.9,0.9}{
            \makecell[{{p{\linewidth}}}]{
                \texttt{\tiny{[P1$\langle$GM]}}
                \texttt{Game Violation: ParseError: Response is not valid JSON.} \\
\texttt{Output ONLY:} \\
\texttt{One JSON with confidence\_score, current\_mode, input\_rationale, and hypothesis keys ...} \\
\texttt{OR} \\
\texttt{SOLVE: ```python ... ```} \\
            }
        }
    }
     \\ \\

    \theutterance \stepcounter{utterance}  
    & \multicolumn{4}{p{0.6\linewidth}}{
        \cellcolor[rgb]{0.561,0.737,0.561}{
            \makecell[{{p{\linewidth}}}]{
                \texttt{\tiny{[P1$\rangle$GM]}}
                \texttt{\{"current\_mode": "probing", "input\_rationale": "I have four examples: f({-}10) = {-}11, f(0) = 0, f({-}2) = {-}4, f({-}99) = {-}198. Let's analyze the ratios: f({-}10)/{-}10 = 1.1, f({-}2)/{-}2 = 2, f({-}99)/{-}99 = 2. It seems for x <= {-}10, f(x) = 2x? No, f({-}10) = {-}20 != {-}11. Wait, f({-}99) = {-}198 = 2 * {-}99. f({-}2) = {-}4 = 2 * {-}2. f(0) = 0 = 2 * 0. But f({-}10) = {-}11 != 2 * {-}10. This is inconsistent. Let's re{-}read f({-}10) = {-}11. Maybe it's not 2x. Let's look at f(x) = x + floor(x/10) again. f({-}10) = {-}10 + ({-}1) = {-}11. f({-}99) = {-}99 + ({-}10) = {-}109 != {-}198. So that's wrong. What if f(x) = 2x for x} \\
            }
        }
    }
    & & \\ \\

    \theutterance \stepcounter{utterance}  
    & & & \multicolumn{4}{p{0.6\linewidth}}{
        \cellcolor[rgb]{0.9,0.9,0.9}{
            \makecell[{{p{\linewidth}}}]{
                \texttt{\tiny{[P1$\langle$GM]}}
                \texttt{Game Violation: ParseError: Response is not valid JSON.} \\
\texttt{Output ONLY:} \\
\texttt{One JSON with confidence\_score, current\_mode, input\_rationale, and hypothesis keys ...} \\
\texttt{OR} \\
\texttt{SOLVE: ```python ... ```} \\
            }
        }
    }
     \\ \\

    \theutterance \stepcounter{utterance}  
    & \multicolumn{4}{p{0.6\linewidth}}{
        \cellcolor[rgb]{0.561,0.737,0.561}{
            \makecell[{{p{\linewidth}}}]{
                \texttt{\tiny{[P1$\rangle$GM]}}
                \texttt{\{"current\_mode": "probing", "input\_rationale": "I have four examples: f({-}10) = {-}11, f(0) = 0, f({-}2) = {-}4, f({-}99) = {-}198. Let's look at the pattern again. f(0)=0, f({-}2)={-}4, f({-}99)={-}198. These suggest f(x) = 2*x for x != {-}10? But f({-}10) = {-}11, not {-}20. This is a specific exception or a different formula. Let's try f(x) = 2*x for x < 0, except when x is a multiple of 10? No, {-}99 is not a multiple of 10. Let's try f(x) = x + floor(x/10) for x <= {-}10? f({-}10) = {-}11. f({-}99) = {-}99 + ({-}10) = {-}109 != {-}198. Let's try f(x) = 2*x for x < {-}10? f({-}99) = {-}198. f({-}10) = {-}11. f({-}2) = {-}4. This suggests a piecewise function. Let's test x={-}11 to see if f({-}11) = {-}22.",} \\
            }
        }
    }
    & & \\ \\

    \theutterance \stepcounter{utterance}  
    & & & \multicolumn{4}{p{0.6\linewidth}}{
        \cellcolor[rgb]{0.9,0.9,0.9}{
            \makecell[{{p{\linewidth}}}]{
                \texttt{\tiny{[P1$\langle$GM]}}
                \texttt{Game Violation: ParseError: Response is not valid JSON.} \\
\texttt{Output ONLY:} \\
\texttt{One JSON with confidence\_score, current\_mode, input\_rationale, and hypothesis keys ...} \\
\texttt{OR} \\
\texttt{SOLVE: ```python ... ```} \\
            }
        }
    }
     \\ \\

    \theutterance \stepcounter{utterance}  
    & \multicolumn{4}{p{0.6\linewidth}}{
        \cellcolor[rgb]{0.561,0.737,0.561}{
            \makecell[{{p{\linewidth}}}]{
                \texttt{\tiny{[P1$\rangle$GM]}}
                \texttt{\{"current\_mode": "probing", "input\_rationale": "I have four examples: f({-}10) = {-}11, f(0) = 0, f({-}2) = {-}4, f({-}99) = {-}198. The pattern f(x) = 2*x holds for {-}2 and {-}99, but fails for {-}10. f({-}10) = {-}11 is an outlier. Let's test x={-}11 to see if it follows the 2*x pattern or the {-}11 pattern.", "hypothesis": "unknown", "confidence\_score": 40\}} \\
            }
        }
    }
    & & \\ \\

    \theutterance \stepcounter{utterance}  
    & & & \multicolumn{4}{p{0.6\linewidth}}{
        \cellcolor[rgb]{0.9,0.9,0.9}{
            \makecell[{{p{\linewidth}}}]{
                \texttt{\tiny{[P1$\langle$GM]}}
                \texttt{OUTPUT: (99, 0)} \\
            }
        }
    }
     \\ \\

    \theutterance \stepcounter{utterance}  
    & \multicolumn{4}{p{0.6\linewidth}}{
        \cellcolor[rgb]{0.561,0.737,0.561}{
            \makecell[{{p{\linewidth}}}]{
                \texttt{\tiny{[P1$\rangle$GM]}}
                \texttt{\{"current\_mode": "probing", "input\_rationale": "I have five examples: f({-}10) = {-}11, f(0) = 0, f({-}2) = {-}4, f({-}99) = {-}198, f(99) = 0. The positive input 99 maps to 0. This suggests that for positive x, f(x) might be 0. Let's check if f(x) = 0 for x > 0. Also, f(0) = 0. So for x >= 0, f(x) = 0. For negative x, we have f({-}2) = {-}4, f({-}99) = {-}198, which is 2*x. But f({-}10) = {-}11, which is not 2*({-}10) = {-}20. This is strange. Let's re{-}examine f({-}10) = {-}11. Maybe it's not 2*x for all negatives. Let's test x={-}1 to see if f({-}1) = {-}2.", "hypothesis": "unknown", "confidence\_score": 50\}} \\
            }
        }
    }
    & & \\ \\

    \theutterance \stepcounter{utterance}  
    & & & \multicolumn{4}{p{0.6\linewidth}}{
        \cellcolor[rgb]{0.9,0.9,0.9}{
            \makecell[{{p{\linewidth}}}]{
                \texttt{\tiny{[P1$\langle$GM]}}
                \texttt{OUTPUT: (10, 9)} \\
            }
        }
    }
     \\ \\

    \theutterance \stepcounter{utterance}  
    & \multicolumn{4}{p{0.6\linewidth}}{
        \cellcolor[rgb]{0.561,0.737,0.561}{
            \makecell[{{p{\linewidth}}}]{
                \texttt{\tiny{[P1$\rangle$GM]}}
                \texttt{\{"current\_mode": "probing", "input\_rationale": "I have six examples: f({-}10) = {-}11, f(0) = 0, f({-}2) = {-}4, f({-}99) = {-}198, f(99) = 0, f(10) = 9. Let's analyze the positive side: f(0)=0, f(10)=9, f(99)=0. This is not simply 0 for all positive x. f(10)=9 suggests f(x) = x{-}1 for x=10. But f(99)=0 suggests something else. Let's look at the negative side: f({-}2)={-}4, f({-}10)={-}11, f({-}99)={-}198. f({-}2) = 2*({-}2), f({-}99) = 2*({-}99), but f({-}10) = {-}11. This is very inconsistent. Let's reconsider the whole set. Maybe the function is related to the number of digits or some other property. Or maybe I misread the outputs. Let's assume the outputs are correct. f(10)=9, f(99)=0. This looks like f(x) = x {-} 1 for x=10, but f(99)=0. What if f(x) = x {-}} \\
            }
        }
    }
    & & \\ \\

    \theutterance \stepcounter{utterance}  
    & & & \multicolumn{4}{p{0.6\linewidth}}{
        \cellcolor[rgb]{0.9,0.9,0.9}{
            \makecell[{{p{\linewidth}}}]{
                \texttt{\tiny{[P1$\langle$GM]}}
                \texttt{Game Violation: ParseError: Response is not valid JSON.} \\
\texttt{Output ONLY:} \\
\texttt{One JSON with confidence\_score, current\_mode, input\_rationale, and hypothesis keys ...} \\
\texttt{OR} \\
\texttt{SOLVE: ```python ... ```} \\
            }
        }
    }
     \\ \\

    \theutterance \stepcounter{utterance}  
    & \multicolumn{4}{p{0.6\linewidth}}{
        \cellcolor[rgb]{0.561,0.737,0.561}{
            \makecell[{{p{\linewidth}}}]{
                \texttt{\tiny{[P1$\rangle$GM]}}
                \texttt{\{"current\_mode": "probing", "input\_rationale": "I have six examples: f({-}10) = {-}11, f(0) = 0, f({-}2) = {-}4, f({-}99) = {-}198, f(99) = 0, f(10) = 9. Let's look at the positive side again: f(0)=0, f(10)=9, f(99)=0. This is weird. f(10)=9 is x{-}1. f(99)=0 is not x{-}1. Let's look at the negative side: f({-}2)={-}4 (2x), f({-}10)={-}11 (x{-}1), f({-}99)={-}198 (2x). It seems like for multiples of 10, it's x{-}1, and for others, it's 2x? Let's test x=11 to see if f(11) = 22 or 10.", "hypothesis": "unknown", "confidence\_score": 40\}} \\
            }
        }
    }
    & & \\ \\

    \theutterance \stepcounter{utterance}  
    & & & \multicolumn{4}{p{0.6\linewidth}}{
        \cellcolor[rgb]{0.9,0.9,0.9}{
            \makecell[{{p{\linewidth}}}]{
                \texttt{\tiny{[P1$\langle$GM]}}
                \texttt{OUTPUT: (1000, 999)} \\
            }
        }
    }
     \\ \\

    \theutterance \stepcounter{utterance}  
    & \multicolumn{4}{p{0.6\linewidth}}{
        \cellcolor[rgb]{0.561,0.737,0.561}{
            \makecell[{{p{\linewidth}}}]{
                \texttt{\tiny{[P1$\rangle$GM]}}
                \texttt{\{"current\_mode": "probing", "input\_rationale": "I have seven examples: f({-}10) = {-}11, f(0) = 0, f({-}2) = {-}4, f({-}99) = {-}198, f(99) = 0, f(10) = 9, f(1000) = 999. Let's analyze the pattern. For positive x: f(10)=9 (x{-}1), f(1000)=999 (x{-}1), f(99)=0. This is inconsistent. f(99)=0 is the outlier. Let's look at negative x: f({-}2)={-}4 (2x), f({-}10)={-}11 (x{-}1), f({-}99)={-}198 (2x). It seems like for multiples of 10, f(x) = x{-}1. For non{-}multiples of 10, f(x) = 2x? Let's check f(99). 99 is not a multiple of 10. If the rule was 2x, f(99) should be 198, but it is 0. So that's not it. Let's look at f(99)=0 again. Maybe it's related to the sum of digits? 9+} \\
            }
        }
    }
    & & \\ \\

    \theutterance \stepcounter{utterance}  
    & & & \multicolumn{4}{p{0.6\linewidth}}{
        \cellcolor[rgb]{0.9,0.9,0.9}{
            \makecell[{{p{\linewidth}}}]{
                \texttt{\tiny{[P1$\langle$GM]}}
                \texttt{Game Violation: ParseError: Response is not valid JSON.} \\
\texttt{Output ONLY:} \\
\texttt{One JSON with confidence\_score, current\_mode, input\_rationale, and hypothesis keys ...} \\
\texttt{OR} \\
\texttt{SOLVE: ```python ... ```} \\
            }
        }
    }
     \\ \\

    \theutterance \stepcounter{utterance}  
    & \multicolumn{4}{p{0.6\linewidth}}{
        \cellcolor[rgb]{0.561,0.737,0.561}{
            \makecell[{{p{\linewidth}}}]{
                \texttt{\tiny{[P1$\rangle$GM]}}
                \texttt{\{"current\_mode": "probing", "input\_rationale": "I have seven examples: f({-}10) = {-}11, f(0) = 0, f({-}2) = {-}4, f({-}99) = {-}198, f(99) = 0, f(10) = 9, f(1000) = 999. Let's look at the pattern for multiples of 10: f({-}10)={-}11, f(10)=9, f(1000)=999. It seems f(10k) = 10k {-} 1 for k>0, and f({-}10k) = {-}10k {-} 1 for k>0. So f(10k) = 10k {-} 1. For non{-}multiples of 10: f({-}2)={-}4, f({-}99)={-}198, f(99)=0. f({-}2) = 2*({-}2), f({-}99) = 2*({-}99). But f(99) = 0. This is strange. Let's test x=1 to see if f(1) = 2 or 0.", "hypothesis": "unknown", "confidence\_score": 40\}} \\
            }
        }
    }
    & & \\ \\

    \theutterance \stepcounter{utterance}  
    & & & \multicolumn{2}{p{0.3\linewidth}}{
        \cellcolor[rgb]{0.9,0.9,0.9}{
            \makecell[{{p{\linewidth}}}]{
                \texttt{\tiny{[GM$|$GM]}}
                \texttt{Game Verdict: LOSS (turn limit reached)} \\
            }
        }
    }
    & & \\ \\

\end{supertabular}
}

\end{document}
